\section{引言}
\label{sec:intro}

高血压是最常见的慢性心血管疾病风险因素，长期控制不佳会显著增加脑卒中与心力衰竭风险\cite{kario2024who}。但袖带式设备依赖间歇充气，难以满足长时程监测需求，因此连续、无感的血压监测具有重大的临床价值。

光电容积脉搏波通过光学方式观测皮下微血管的血容量变化，可集成于智能手表、指环与耳戴设备，具备低成本、低功耗、易佩戴的优势。其波形形态与动脉顺应性、外周阻力、脉搏波传导时间等因素相关，因而与收缩压、舒张压存在统计关联\cite{charlton2022wearable,tang2026sifpbpnet}；但这种关联受血管张力、体温、体位与个体差异调制，并不构成确定的物理方程，估计问题最终归结为映射学习，精度取决于表征形式、监督标签与验证协议三者的匹配。本文以如何由 PPG 估计血压为主线，围绕表征、监督与验证三个环节进行综述。

\noindent 本文的主要内容可概括为：梳理 PPG 的生理学基础与信号特征，说明其与血压的统计关联及局限；按手工特征与传统回归、端到端深度网络、预训练并微调三条路线，比较各类方法在表征、监督与部署代价上的取舍，并讨论个体化校准与多模态融合；汇总常用数据集与评价标准，分析验证协议对结论可靠性的影响，并分析人群差异、校准与标签稀缺等问题。

本文组织结构如下：\cref{sec:physio}介绍生理基础与波形特征；\cref{sec:techniques}比较三类估计方法；\cref{sec:datasets}介绍数据集与评价标准；\cref{sec:challenges}总结挑战与未来方向；\cref{sec:conclusion}给出结论。